\begin{minipage}{0.6\linewidth}
Une échelle uniforme de \qty{\v{M}}{kg} et de \qty{\v{L}}{m} est appuyée contre un mur vertical sans frottement et fait un angle de \qty{\v{th}}{\degree} avec le sol horizontal. Une personne de \qty{\v{m}}{kg} est montée de \qty{\v{s}}{m} le long de l'échelle. L'échelle ne glisse pas.
\end{minipage}\hfill
\begin{minipage}{0.35\linewidth}
\centering
\begin{tikzpicture}[scale=0.8]
  \fill[pattern=north east lines] (-0.5,-0.3) rectangle (3.1,0); \draw[thick] (-0.5,0) -- (3.1,0);
  \fill[pattern=north east lines] (2.5,0) rectangle (2.8,4); \draw[thick] (2.5,0) -- (2.5,4);
  \draw[line width=2pt] (0.9,0) -- (2.5,3.6);
  \draw (1.5,0) arc (0:66:0.6); \node at (1.8,0.35) {$\theta$};
\end{tikzpicture}
\end{minipage}

\begin{enumerate}[label=\alph*)]
    \item Déterminez la force normale exercée par le mur sur l'échelle.
    \item Déterminez la force normale et la force de frottement exercées par le sol sur l'échelle.
    \item Déterminez le coefficient de frottement statique minimal entre l'échelle et le sol.
\end{enumerate}
